\section{Recovering the two polar lines in the full completion}
\label{sec:trace-polar-completion}

The principal finite subspace with two vanishing moments is not
the full subspace in C99.  The following construction specifies
the Hilbert completion and proves the contribution of its two
quotient lines.  Its finite families are the actual even prolate
families of AC-T42--T44.  It does not identify every possible
adelic norm cutoff with that chosen family.

For an even Schwartz function on $\mathbb R$, put
\[
 I\psi=\int_0^\infty\psi(x)\,dx,\qquad
 E_\zeta\psi(t)=e^{t/2}\sum_{n\geq1}\psi(ne^t)
                         -I\psi\,e^{-t/2},\qquad
 J\psi=(E_\zeta\psi,\psi(0),I\psi).
 \tag{PC.1}
\]
Write $\mathcal H=L^2(\mathbb R,dt)\oplus\mathbb C^2$ with
the product Hilbert norm.  Let $\mathcal P_\Lambda$ be the
retained, time-supported even prolate space, with the same
threshold as AC-T42, and define
\[
 \mathcal B_\Lambda=J\mathcal P_\Lambda,\qquad
 \mathcal B_{\Lambda,0}
  =J\{\psi\in\mathcal P_\Lambda:\psi(0)=I\psi=0\}.
 \tag{PC.2}
\]
The corrected transform on these piecewise smooth functions
is the extension established in AC-T42.

\begin{theorem}[AC-T45: the full principal completion and polar trace]
\label{AC-T45}
The separated completion of the even Schwartz space for
$\|J\psi\|_{\mathcal H}$ is $\mathcal H$.  The completion
of its two-moment-zero subspace is
$\mathcal H_0=L^2(\mathbb R)\oplus0$.
Normalized dilation
$R(a)\psi(x)=a^{-1/2}\psi(x/a)$ acts in this completion as
\[
 R(a)=V(a)\oplus a^{-1/2}\oplus a^{1/2},\qquad
       (V(a)F)(t)=F(t-\log a).
 \tag{PC.3}
\]
Let $P_\Lambda,P_{\Lambda,0}$ be the orthogonal projections
onto the two spaces in (PC.2).  Then
\[
 P_{\Lambda,0}\longrightarrow P_{\mathcal H_0}
       \quad\hbox{strongly},\qquad
 P_\Lambda-P_{\Lambda,0}\longrightarrow P_{\mathbb C^2}
       \quad\hbox{in trace norm}.
 \tag{PC.4}
\]
The second assertion, including the rank two, holds for all
sufficiently large $\Lambda$ and requires no RH assumption.
Consequently, for every $h\in C_c^\infty(\mathbb R_+^\times)$,
\[
 \begin{split}
 \operatorname{Tr}(P_\Lambda R(h))
   &=\operatorname{Tr}(P_{\Lambda,0}R(h))
       +\int_0^\infty h(a)(a^{-1/2}+a^{1/2})\,d^*a+o(1).
 \end{split}
 \tag{PC.5}
\]
Under $\mathsf A$, combining (PC.5) with AC-T44 gives
\[
 \begin{split}
 \operatorname{Tr}(P_\Lambda R(h))
  ={}&2\log\Lambda\,h(1)
     -\sum_\gamma m_\gamma\widehat h_+(\gamma)\\
    &+\int_0^\infty h(a)(a^{-1/2}+a^{1/2})\,d^*a+o(1).
 \end{split}
 \tag{PC.6}
\]
Here the zero sum uses the same Fourier convention as AC-T44;
for the zeta sector it is unchanged by reversing all ordinates.
For the unnormalized action $U(a)\psi(x)=\psi(x/a)$,
the two additional terms are exactly
$\int h(a)\,d^*a+\int a h(a)\,d^*a$.
\end{theorem}
\begin{proof}
\proofstep{1. The corrected transform and the extension.}
Poisson summation for an even Schwartz function gives
\[
 E_\zeta\psi(t)=E_\zeta\widehat\psi(-t),\qquad
 I\widehat\psi=\psi(0)/2.
 \tag{PC.7}
\]
The convention is $\widehat\psi(y)=\int_\mathbb R
\psi(x)e^{-2\pi ixy}\,dx$.
At $+\infty$, the sum in (PC.1) decreases faster than any
exponential of fixed rate, while its remaining term is
$-I\psi e^{-t/2}$.
At $-\infty$, (PC.7) gives the corresponding leading term
$-\psi(0)e^{t/2}/2$.
In particular $E_\zeta\psi\in L^2(\mathbb R)$ for every
even Schwartz function; neither moment has been discarded.
Changing variables in (PC.1) proves (PC.3) directly.
The last two coordinates of $J$ are surjective: an even
smooth bump supported away from zero has nonzero integral
and zero value at zero, while a bump equal to one at zero
can have its integral removed by a multiple of the former.

\proofstep{2. Density of the zero-moment image.}
Choose a nonzero even smooth compactly supported $\psi$
which vanishes near zero and has $I\psi=0$.
Such a function is obtained by subtracting suitable multiples
of two disjoint positive bumps on the positive half-line
and reflecting them evenly.
Its Mellin transform
$m(s)=\int_0^\infty\psi(x)x^{s-1}\,dx$ is entire and
not identically zero, by Fourier uniqueness after $x=e^t$.
The Mellin identity for the corrected theta transform is
\[
        \widehat{E_\zeta\psi}(\tau)
           =\zeta(1/2+i\tau)m(1/2+i\tau)
 \tag{PC.8}
\]
up to the harmless simultaneous sign convention for Fourier
transforms.  This identity is AC-T42's Mellin formula;
it also follows by subtracting the integral term in the
usual Mellin summation identity on $0<\Re s<1$.
The product in (PC.8) is nonzero almost everywhere on the
critical line, since each of its factors is a nonzero
holomorphic function in a neighborhood of that line and
its zeros are discrete.  This uses no zero-location assertion.

If $F\in L^2(\mathbb R)$ is orthogonal to every translate
of $E_\zeta\psi$, then the Fourier transform of the
$L^1$ function $\widehat F\,\overline{\widehat{E_\zeta\psi}}$
vanishes identically.  Fourier uniqueness implies that
this product is zero almost everywhere, hence $F=0$.
Every such translate is the image of a normalized dilation
of $\psi$, still smooth, compactly supported, vanishing
near zero and of integral zero.  Thus these images are
dense in $L^2(\mathbb R)$.

The norm on the zero-moment subspace is exactly its image
norm.  The preceding density proves its completion is
$\mathcal H_0$.  Given $(F,z)\in L^2\oplus\mathbb C^2$,
choose a fixed smooth compactly supported $f_z$ with
$(f_z(0),If_z)=z$.  Approximate $F-E_\zeta f_z$ by
zero-moment images and add $Jf_z$.
This proves that the completion of the full space is
all of $\mathcal H$ and realizes the exact sequence
$0\to\mathcal H_0\to\mathcal H\to\mathbb C^2\to0$.

\proofstep{3. The actual finite subspaces are dense.}
The fixed-profile conclusion of AC-T42 states that for every
even smooth compactly supported zero-moment $f$ there are
$f_\Lambda\in\mathcal P_\Lambda$ with the same two zero
moments and
$\|E_\zeta(f_\Lambda-f)\|_2\to0$.
The proof there uses the general graph estimate with the
Schwartz Fourier tail, and corrects the two moments using
two retained profiles with an invertible constraint matrix.
In particular it applies to each of the fixed dilation
profiles in step 2; it does not require the spaces
$\mathcal P_\Lambda$ to be nested.
The distance of each member of a dense set in
$\mathcal H_0$ to $\mathcal B_{\Lambda,0}$ tends to zero.
Contractivity of orthogonal projections extends this
assertion to every member of $\mathcal H_0$.
On $\mathbb C^2$ the projections are zero, proving the
first part of (PC.4).

\proofstep{4. An exact rank-two graph.}
Choose fixed even smooth compactly supported $f_1,f_2$
with $(f_j(0),If_j)=e_j$, where $e_1,e_2$ is the standard
basis of $\mathbb C^2$.
The unconstrained fixed-profile statement of AC-T42 gives
$\widetilde f_{\Lambda,j}\in\mathcal P_\Lambda$ such that
\[
 E_\zeta\widetilde f_{\Lambda,j}\to E_\zeta f_j,
 \qquad
 (\widetilde f_{\Lambda,j}(0),I\widetilde f_{\Lambda,j})
                         \to e_j.
 \tag{PC.9}
\]
Its proof controls evaluation by a local Sobolev inequality
and the integral by the time-support Cauchy--Schwarz bound,
so that (PC.9) does not infer point convergence from
$L^2$ convergence alone.
The $2\times2$ matrix of their moments tends to the identity.
For large $\Lambda$ multiply this pair by its inverse, to
obtain $f_{\Lambda,j}$ with moments exactly $e_j$ and
$E_\zeta f_{\Lambda,j}\to E_\zeta f_j$.
Thus the moment map on $\mathcal P_\Lambda$ is surjective.

Write $Q_{\Lambda,0}$ for the projection in $L^2$ corresponding
to $P_{\Lambda,0}$.  The orthogonal complement of
$\mathcal B_{\Lambda,0}$ inside $\mathcal B_\Lambda$ is
exactly the two-dimensional span of
\[
 w_{\Lambda,j}
   =((1-Q_{\Lambda,0})E_\zeta f_{\Lambda,j},e_j),
                              \qquad j=1,2.
 \tag{PC.10}
\]
Indeed subtracting $P_{\Lambda,0}Jf_{\Lambda,j}$ leaves
these vectors in that complement, and any vector in
$\mathcal B_\Lambda$ with zero moments already lies in
$\mathcal B_{\Lambda,0}$.
Strong convergence from step 3 and (PC.9) show
$w_{\Lambda,j}\to(0,e_j)$ in norm.
The Gram matrix of the two vectors tends to the identity.
The formula $W_\Lambda(W_\Lambda^*W_\Lambda)^{-1}W_\Lambda^*$
for their projection now proves operator-norm convergence
to $P_{\mathbb C^2}$.  The difference has rank at most four,
so its trace norm is at most four times its operator norm.
This proves the second assertion of (PC.4).

\proofstep{5. The trace and its source meaning.}
For compactly supported $h$, $R(h)$ is bounded, including
on the two nonunitary finite-dimensional summands.
The trace-norm assertion gives
\[
 \left|\operatorname{Tr}\bigl(
 (P_\Lambda-P_{\Lambda,0}-P_{\mathbb C^2})R(h)\bigr)\right|
 \leq\|P_\Lambda-P_{\Lambda,0}-P_{\mathbb C^2}\|_1\|R(h)\|
 \longrightarrow0.
 \tag{PC.11}
\]
The trace of $R(h)$ on $\mathbb C^2$ is precisely the
integral in (PC.5).  AC-T44 supplies the trace on
$\mathcal B_{\Lambda,0}$ under $\mathsf A$, proving (PC.6).
Finally $U(a)=a^{1/2}R(a)$, so replacing $h(a)$ by
$a^{1/2}h(a)$ changes the two centered exponents to $0,1$.
These are exactly the quotient characters in C99,
III.(6), (10)--(11), (31)--(34).  In VIII.(21)--(26)
the source uses the zero-moment projection and attributes
the distribution $D$ to its difference from the full
projection.  Equations (PC.4)--(PC.11) prove that difference
for the explicit number-field prolate completion here.
The paragraph after VIII.(33) supplies the source even
archimedean family and the map $\psi\mapsto\psi\otimes1_R$.
No unspecified quotient metric is used in this proof.
\end{proof}

\paragraph{Scope and geometric dictionary.}
The two finite lines have characters $1$ and $|\cdot|$ for
$U$, as in the formal $H^0$ and Tate $H^2$ terms of C99's
Lefschetz expression.  This is an equality of characters
of an explicitly constructed representation and an actual
limit of finite projection traces.  It does not construct
an arithmetic surface or assert that its Hilbert completion
is an $\ell$-adic cohomology group.  The exact adelic
fiber-cutoff comparison left in AC-T11 is still distinct
from this now explicit polar-line calculation.
